\chapter{Мера суммы по spin-структурам}

Предыдущий relative-determinant gate показал, что при одинаковых
предварительных весах periodic и anti-periodic ветвей фермионная finite
part ниже для \(\beta=1/2\). Теперь необходимо проверить, содержит ли
текущий parent action саму сумму, относительно которой это ранжирование
становится динамическим выбором.

\section{Что определено текущим parent triple}

Для каждой spin-структуры \(\mathfrak s\) почти-коммутативный lift имеет вид
\begin{align}
 \mathcal H_{\mathfrak s}
 &=L^2(K,S_{\mathfrak s})\otimes H_8,
 \\
 \mathcal D_{\mathfrak s,\rho}
 &=D_{K,\mathfrak s,\rho}\otimes\mathbf1
 +\gamma_K\otimes\chi D_F^*.
\end{align}
Это определяет отдельный partition function
\begin{equation}
 Z_{\mathfrak s,\rho}
 =\int_{\mathcal F_{\mathfrak s,\rho}}
 e^{-S_{\mathfrak s,\rho}}\,D\mathcal F.
\end{equation}
BRST/BV complex действует на полях внутри фиксированного spinor bundle.
Ни одно его поле не переводит \(S_{\mathfrak s}\) в
\(S_{\mathfrak s'}\), а finite algebra не содержит spin-transition edge.

Следовательно, выражение
\begin{equation}
 Z_{\mathrm{sum}}
 =\sum_{\mathfrak s}w_{\mathfrak s}Z_{\mathfrak s}
 \label{eq:spin-sum-partition}
\end{equation}
не является следствием уже заданного functional integral. Оно определяет
дополнительную теорию и требует меры \(w_{\mathfrak s}\).

\section{Может ли равномерная мера считаться канонической}

Spin-структуры образуют torsor над
\begin{equation}
 H^1(K;\mathbb Z_2)\cong\mathbb Z_2^2.
\end{equation}
Если заранее постулировать физическое действие всей torsor translation
group, единственная нормированная инвариантная counting measure равна
\begin{equation}
 w_{\mathfrak s}=\frac14.
\end{equation}
Но наличие torsor не означает, что его абстрактные трансляции являются
gauge symmetry parent action. Текущая алгебра автоморфизмов, Dirac operator
и BV complex такого действия не выводят. Поэтому равномерная мера является
естественным дополнительным постулатом, но не theorem текущей версии.

Неравномерные фазы или знаки также не фиксированы. Их можно получить только
после добавления, например, \(\mathbb Z_2\)-topological gauge sector,
spin-TQFT, determinant-line weight или явно объявленного filling/cobordism
prescription. Каждый такой вариант меняет parent theory и обязан заново
пройти anomaly, reality и measure gates.

\begin{theorem}[Spin-sum specification gap]
Текущий base-\(K\) parent action определяет четыре отдельных теории на
фиксированных spin-структурах, но не их сумму. Поэтому более низкое значение
relative determinant при \(\beta=1/2\) является scheme-independent
ранжированием ветвей, а не динамическим выбором внутри уже построенного
partition function.
\end{theorem}

\section{Конфликт с прежней periodic постановкой}

В старой RPFT-ветви periodic condition на \(S^1\) фиксировалось из чтения
окружности как пространственной, а anti-periodic condition называлось
термальным. Это различение слишком сильное. Anti-periodicity является
обязательной KMS-конвенцией, когда окружность представляет евклидово время;
на пространственной окружности periodic и anti-periodic условия являются
двумя допустимыми spin-структурами.

Следовательно, spatial interpretation сама по себе не выводит periodic
branch. Старый выбор остаётся допустимым model input, но не может
одновременно использоваться вместе с утверждением, что determinant
динамически выбрал anti-periodic branch, пока не задана сумма
\eqref{eq:spin-sum-partition}.

\begin{nogocriterion}[Запрет подмены background selection]
Сравнение \(Z_{\rm AP}/Z_{\rm P}\) не разрешает заменить заранее выбранную
periodic spin-структуру на anti-periodic. Для такого перехода требуется
либо общий spin-sum partition function с выведенными весами, либо
динамический topological field, связывающий sectors.
\end{nogocriterion}

\section{Замороженный статус}

Итог имеет двухуровневый вид:
\begin{equation}
 \boxed{
 \begin{gathered}
 \Delta\Gamma_f(1/2)<\Delta\Gamma_f(0)
 \quad\text{--- доказанное relative ranking},
 \\
 \beta_{S^1}=1/2
 \quad\text{--- не выведенный dynamic selection}.
 \end{gathered}}
\end{equation}
Минимальное продолжение без смены теории состоит в явной заморозке одной
spin-структуры как background train input. Альтернативное продолжение ---
построение новой \(\mathbb Z_2\)/spin-TQFT measure; оно относится уже к
расширенной версии parent action.

\begin{protocol}[Следующий допустимый выбор]
До вычисления полного численного zeta-потенциала необходимо выбрать один из
двух путей:
\begin{enumerate}
 \item сохранить legacy periodic branch как явный spatial-background input;
 \item принять anti-periodic branch как новый явный input, мотивированный
 relative determinant, но не называть его динамически выведенным.
\end{enumerate}
Суммирование с равными весами допустимо только после отдельного объявления
spin-torsor symmetry как части parent theory.
\end{protocol}

Машинный статус сохранён в
\path{s2t_v4_spin_sum_measure_gate_results.json}.